\chapter{序論}
\label{chap:introduction}

日常生活の記録は、歩数のように測定しやすい量から、端末操作、予定、周辺環境、本人の主観を関連付けるライフログへと対象を広げてきた。しかし、記録された活動量と、本人が経験する生活状態は同一ではない。同じ場所で同じ程度の入力をしていても、仕事、学習、私的な連絡など、異なる活動が含まれる。固定された活動名だけでは表しきれない個人の生活のまとまりを、観測と本人確認の両方から捉える必要がある。

本研究は、研究用キーボードと任意同意の生活情報を用い、拠点、個人固有の行動パターン、現在・日次の5尺度を区別して表現する三階層生活状態モデルを構築する。活動カテゴリをあらかじめ24種類に固定して当てることを主目的とせず、未知の行動候補を発見し、本人の確認・命名・訂正と結び付ける。その成立条件と限界を、探索的予備研究で検討する。

\section{身近な活動量計測から生活の文脈へ}

歩数は身近な身体活動を継続的に記録し、日ごとの違いを振り返る手段になる。対象集団や測定条件を考慮して解釈すべき指標でもある\autocite{tudorlocke2011steps}。同じ日歩数であっても、短い歩行を分散して行った日と、一度の外出でまとめて歩いた日では、活動の構成が異なる。機器を携行していない時間の記録不足を、歩かなかったことと同一視することもできない。

本人の過去の記録は、一律の目標とは異なる比較基準を与える。ただし、通常より歩数が少ないことから、疲労や気分の低下を直接断定することはできない。活動量に加え、それが生じた時間、作業、拠点、本人の経験を関連付ける必要がある。

\section{異なる情報源を結び付けるライフログ}

ライフログは、多様な生活記録を時間軸に沿って蓄積し、検索や振り返り、行動理解へ利用する研究領域である。Gurrinらは取得・管理・分析・利用の観点からこの領域を整理している\autocite{gurrin2014lifelogging}。キーボード操作は入力行動、歩数は身体活動、予定は計画された文脈、開閉センサは環境変化をそれぞれ観測する。

それぞれの情報が示す範囲は異なる。予定が存在しても実行されたとは限らず、扉の開閉だけで本人の所在は確定できない。文章の内容も、本人が今実行している作業の正解とは限らない。情報を統合する際には、値、取得条件、時刻、欠測を保持し、観測を意味の確定に置き換えない設計が必要である。

入力操作は高頻度イベント、歩数は累積スナップショット、予定は区間、自己報告は特定の時点を対象とする回答として得られる。さらに通信遅延がある。発生時刻だけを整列しても、予測当時に知り得た情報は再現できない。受信時刻を保持する因果的なデータ構築を、モデルの前提とする。

ライフログの観測だけでは、本人にとっての作業の意味や、発見された候補の区別を確定できない場合がある。著者らが検討した、母音列から日本語候補を生成する省入力方式\autocite{tanaka2026fss}を、本修論では、本人の確認・命名・訂正を受け取る入力接点に関する先行検討として位置付ける。その開発過程では、キーボードを日常行動の取得接点として捉え、さらに観測だけでは分からない意味を本人に確認する窓口として利用する可能性に着目した。

この先行検討では、打鍵削減の可能性と母音変換等に伴う認知的負担が併存した。本修論では、この入力接点への問題関心を背景に、異種観測からの生活パターン発見と本人による確認・命名を結び付け、取得の継続性と回答負担も検討する。先行方式とその予備評価は第\ref{chap:related-work}章2.2.3節で述べる。

\section{固定分類と個人固有のパターンの隔たり}

「勉強・研究」という一つのカテゴリでも、論文執筆、プログラミング、実験など、本人にとって区別したい作業がある。一方、固定分類と別のTask分類を並行して用意すると、どちらが行動の定義を担うかが曖昧になる。固定分類器の出力と個人固有の行動発見を、同じ目的の二つの推定として重複させないことが課題となる。

時間尺度の違いもある。入力間隔や短い中断は数分単位で変化するが、日内の作業の流れや曜日の傾向は、より長い系列を必要とする。単一窓への集約だけでは細かな変化と長期文脈を両立しにくいため、時間尺度に応じた表現を段階的に構成する。

ラベルを減らすことは、意味を機械が自動で理解することではない。自宅でしか行わない作業では、場所の特徴と作業の特徴が共起する。観測だけでは両者を分離できない場合がある。候補を提示し、同じ場所か、同じ作業かを別々に本人が確認できる仕組みが必要である。

自己報告には負担と偏りが伴う。EMAは日常場面で主観状態を得る手段であるが、提示頻度、回答遅延、未回答も測定過程の一部となる\autocite{shiffman2008ema}。忙しい時点ほど回答しにくい場合、得られた教師は生活全体の無作為標本ではない。推定誤差だけでなく、取得の継続性と回答負担を評価する。

\section{本研究の目的と提案構成}

目的は、異種観測から個人の生活の反復を表現し、拠点・行動パターン・主観尺度を区別したまま、本人の確認と修正に接続するモデルを構築することである。研究用キーボードは観測手段であり、文章の分類そのものを研究目的としない。文章特徴は研究サーバ上のローカルモデルで生成し、原文を外部AIサービスへ送信しない。

共有エンコーダにはMicro・Meso・Macroの三階層LSTMを用いる。Microは5分特徴の区間内系列、Mesoは日内の区間系列、Macroは当週の日表現と先行28日の日別観測を処理する。LSTMの記憶機構を異なる時間単位へ適用し、上位文脈を現在・日次の表現へ融合する\autocite{hochreiter1997lstm}。

共有表現からPlaceとBehaviorの埋め込み、および現在・日次の5尺度を出力する。行動埋め込みには閾値型の逐次プロトタイプ法を適用し、カテゴリ数を事前に固定せず候補を作る。自己教師ありの再構成・将来予測と、本人確認による関係学習・尺度教師を同時に用いる。拠点と行動ごとに別の三階層モデルを動かすものではない。

24活動の既存回答は削除せず、学習後の意味注釈と従来分類器の比較に用いる。提案モデルの固定出力、交差エントロピー、入力区間境界には使わない。候補の識別子（ID）、本人が付けた名前、モデルによる推定、本人確認による確定をそれぞれ区別する。

現行実装はバッチ再学習と凍結モデル・マップによる推論である。未知観測からの候補生成は再学習時に行い、オンライン重み更新や自動的な意味の獲得を実装済みとは主張しない。新しい状態に基づく歩数基準と運動・休憩提案は別途検証する。状態を発見することと、有効な支援へつながることを分けて研究する。

\section{評価の範囲と論文構成}

最大5名・各14日間の探索的予備研究を計画し、取得の実現可能性と回答負担を主評価とする。拠点・行動のID一致、未知と候補の扱い、場所と行動の分離、現在・日次尺度の誤差を補助的に調べる。合成テストの成功を実参加者での有効性と混同せず、成立しなかった条件も記録する。

母音省入力の先行検討で打鍵削減と認知的負担が併存したことからも、操作の少なさと継続的な回答のしやすさは同義ではなく、取得の継続性と回答負担を分けて評価する必要がある。

第\ref{chap:related-work}章で関連研究と数理的基盤を整理する。第\ref{chap:observations}章で観測、受信時刻、欠測、本人確認と注釈の役割を定義し、第\ref{chap:model}章で共有三階層表現、学習損失、パターン発見、表示・更新条件を定式化する。第\ref{chap:evaluation}章で評価計画を定め、第\ref{chap:conclusion}章で到達点と未検証事項を整理する。従来の固定分類、独立ContextLSTM、部分プーリングや状態保持の拡張は、提案実装との違いが分かるように扱う。
